%*****************************************
\chapter{Source code overview}\label{app:source-code}
%*****************************************

All experiments in this thesis were run on a Linux laptop with an AMD Ryzen 7 6800HS CPU and 16 GB of RAM. The code for the functional regression experiments was executed with a random seed set to the value 2024 for reproducibility, while the functional testing experiments used a random seed set to 2026. 

The Python code developed for this work is available under a GPLv3 license at the following GitHub repositories: 
\begin{center}
  \begin{tabular}{l}
\url{https://github.com/antcc/rk-bfr-jump},\\
\url{https://github.com/antcc/rk-bt-jump}.
  \end{tabular}
\end{center}

The code is adequately documented and is structured in several directories as follows:

\begin{itemize}
  \item In the \texttt{rkbfr\_jump} and \texttt{rkbt\_jump} folders we find the files responsible for the implementation of our Bayesian models, separated according to the functionality they provide.
  \item The \texttt{reference\_methods} folders contain our implementation of the functional comparison algorithms that were not available through a standard Python library.
  \item The \texttt{results} folders contain plain text files with the execution times and the numerical results, as well as \texttt{.csv} files that facilitate working with them.
  \item At the root folder we have a Python script \texttt{experiments.py} for executing our experiments, which accepts several user-specified parameters (such as the number of iterations or the type of data set).
\end{itemize}

When possible, the code was implemented in a generic way that would allow for easy extensions or derivations. It was also developed with efficiency in mind, so many functions and methods exploit the vectorization capabilities of the \slsc{numpy} and \slsc{scipy} libraries, and are sometimes parallelized using \slsc{numba}. Moreover, since we followed closely the style of the \slsc{scikit-learn} and \slsc{scikit-fda} libraries, our methods are compatible and could be integrated (after some minor tweaking) with both of them. 

Lastly, there are \slsc{Jupyter} notebooks that demonstrate the use of our methods in a more visual way. Inside these notebooks there is a step-by-step guide on how one might execute our algorithms, accompanied by many graphical representations, and offering the possibility of changing multiple parameters to experiment with the code.

%*****************************************
%*****************************************
%*****************************************
%*****************************************
%*****************************************
